
Desátý případ užití je specifikován v tabulce č. \ref{tab:usecase:UC10}.

\begin{UseCaseTable}
  {Příjem a validace pozorování}
  {UC10}
  {F10}
  {Ověření a uložení pozorování pořízeného senzorem, případně odmítnutí neplatného pozorování}
  {Senzor}
  {}
  {Senzor je v systému evidován v aktivním stavu}
  {Pozorování je uloženo do evidence, případně bylo odmítnuto chybovou odpovědí}

\UseCaseScenario{Základní scénář}
\UseCaseStep{senzor}{Odešle systému pozorování (hodnotu, časové razítko a identifikátor veličiny) společně s identifikací zařízení.}
\UseCaseStep{systém}{Ověří, že odesílatel je v systému evidovaným senzorem v aktivním stavu.}
\UseCaseStep{systém}{Ověří, že veličina a jednotka pozorování odpovídají deklaracím v evidenci senzoru.}
\UseCaseStep{systém}{Ověří, že hodnota leží v přípustném rozsahu pro danou veličinu.}
\UseCaseStep{systém}{Uloží pozorování do evidence a vrátí potvrzení o přijetí.}

\UseCaseScenario{Alternativní scénář -- neshoda veličiny nebo jednotky}
\UseCaseStepCustom{A1}{senzor, systém}{shodné s kroky 1--2 základního scénáře.}
\UseCaseStepCustom{A2}{systém}{Zjistí, že veličina nebo jednotka pozorování neodpovídá deklaraci senzoru v evidenci.}
\UseCaseStepCustom{A3}{systém}{Odmítne pozorování chybovou odpovědí popisující povahu rozporu.}

\UseCaseScenario{Alternativní scénář -- hodnota mimo přípustný rozsah}
\UseCaseStepCustom{B1}{senzor, systém}{shodné s kroky 1--3 základního scénáře.}
\UseCaseStepCustom{B2}{systém}{Zjistí, že hodnota leží mimo přípustný rozsah pro danou veličinu.}
\UseCaseStepCustom{B3}{systém}{Odmítne pozorování chybovou odpovědí popisující povahu chyby.}

\UseCaseScenario{Alternativní scénář -- senzor není v aktivním stavu}
\UseCaseStepCustom{C1}{senzor}{shodné s krokem 1 základního scénáře.}
\UseCaseStepCustom{C2}{systém}{Zjistí, že odesílající senzor není v evidenci nebo není v aktivním stavu.}
\UseCaseStepCustom{C3}{systém}{Odmítne pozorování chybovou odpovědí.}

\end{UseCaseTable}
